%!TEX root = ../main.tex
\section{Conclusion}
\label{sec:conc}

We presented an honesty-audited late-fusion pipeline for the ComParE 2017 Cold sub-challenge on URTIC that combines a frozen WavLM-Large backbone with a per-layer-audited prior (A2.5) and a constrained-K late-fusion stage (A5b) admitting handcrafted feature groups by subtractive cold-vs-speaker honesty. The locked configuration A5b K=2 (A2.5 + G4\_gain\_invariant + G5\_modulation) reaches devel\_test UAR $0.7037 \pm 0.0060$ at $N{=}5$, with both speaker-probe substrates passing the 0.0780 ceiling. A 5-seed mean-logit ensemble matches the 2017 hidden-test baseline of 0.710 within measurement noise (0.7090, $\Delta\,-0.001$).

Alongside the controlled positive system, the project closes a three-level de-confounding ladder negatively at every frozen-backbone scope tractable in budget: data-level cross-speaker embedding mixup has an empty $\alpha$-axis operating window (M9); representation-level head-only contrastive with speaker-masked positives is dominated by a simpler cold-CE bottleneck and is purely subtractive when added to CE (M10, M11); gradient-level DANN/MDD adversary at both the projection bottleneck and the un-bottlenecked 4096-d substrate is substrate-resistant — the discriminator memorises individual training-chunk speaker IDs without learning generalising speaker structure (M12, M13, M14).

The seven negative-control disciplines that emerged from this closure (M8 audio-splice self-control, M9 $\alpha$-endpoint controls, M10 bottleneck-confound, M11 subtractive-objective interaction, M12 memorisation-gap in-flight check, M13 disc-ceiling-before-$\lambda$-sweep ordering, M14 substrate noise floor) plus a transferable standalone-UAR-predicts-K-fusion-admission heuristic for honesty-audited late-fusion stacks together form an instrumentation framework for any small-data paralinguistic foundation-model pipeline. We argue this framework, together with the systematic three-level closure pattern and the controlled-positive cumulative-stack ($+0.073$ over leak-corrected baseline), is the load-bearing contribution.

\paragraph{Reproducibility.} All cell sources, results JSONs, and per-rung diagnostic logs are tracked in the project repository under \texttt{model/run.ipynb} and \texttt{results/}. Trained head checkpoints are cached under \path|cache/microsoft_wavlm-large/|; downstream rungs reference checkpoints by seed and rung name (e.g., \path|head_A2grouped_honestprior_seed{42, 123, 7, 999, 31337}.pt|). We document every locked $\beta^*$ and $\tau^*$ per seed, every standalone candidate UAR, and every adversary-health diagnostic in the JSONs to enable straightforward replication.

\paragraph{Future work.} Within the frozen-backbone commitment, two HuBERT-base variants have now been ruled out: mean-pooled across layers (standalone cold-LR UAR $0.5396$, M14 pre-flight failed; Section~\ref{ssec:k3_hubert}) and \emph{with learned layer-weighted softmax} (standalone $0.6217 \pm 0.029$ cleared the M14 threshold but K=3 fusion failed by logit-correlation with the WavLM-A2.5 anchor; Sections~\ref{ssec:k3_hubert}, \ref{ssec:standalone_predictor}). The latter result quantitatively decomposed the architectural advantage into $\sim 71\%$ from the audited LW softmax pooling and $\sim 29\%$ from backbone capacity / pretraining-corpus diversity --- both contribute, refuting the strong form of either hypothesis. Two ensemble-aggregation axes (calibration / stacked weighting and test-time augmentation) have also been ruled out within the architectures tested (Section~\ref{ssec:ensemble_defence}); these closures contributed M15--M17 to the methodology table. The remaining untested variants therefore narrow to: \emph{Multi-FM (recipe-orthogonal candidates only):} (i) HuBERT discrete-token histograms, which capture syllable-rate and utterance-rhythm patterns that pooled-stat representations average out --- a structurally-different feature with much lower per-token speaker-info capacity, and crucially \emph{not subject to the audit-recipe-mirror correlation issue} that closed HuBERT-A2.5 at K=3 (highest-information remaining axis). (ii) HuBERT-large mean-pooled, which has $\sim 3.5\times$ more capacity than HuBERT-base --- a confirmation test of the $\sim 29\%$ capacity contribution we measured on HuBERT-A2.5. \emph{Classical baseline alignment:} the original 2017 ComParE Cold sub-challenge baseline used a 6373-d ComParE-2016 acoustic feature set with an SVM classifier; running this baseline within our speaker-grouped sub-split discipline would resolve the apples-to-oranges devel-vs-hidden-test comparison and produce a strict-alignment number. \emph{Beyond the rejected TTA designs:} M17 narrows future TTA work to perturbations that survive WavLM's input normalisation (additive noise with zero mean) or operate at the pooled-stats stage (SpecAugment-style masking on the $[L, 4D]$ pooled tensor); waveform-level multiplicative or temporal perturbations are now empirically ruled out for FM-substrate ensembling on cold detection. \emph{Outside the frozen-backbone commitment:} a full transformer fine-tune at the layer-weight-open scope (with the M10/M12/M13 control discipline applied) could test whether the de-confounding closure generalises beyond frozen-substrate scopes --- a meaningfully separate research question. The HuBERT discrete-token variant (i) is now the highest-priority untested path because it bypasses the structural correlation that closed HuBERT-A2.5.
